\ea\label{ex:text26-1}
O kaku vagana to dao raara bona vavaaiku na pakupaku raara namana beara hae batara teo sinivi.\\
\gll o	kaku	vagana	to	dao	ra=ara	bona	va-vaa-iku	na	paku{$\sim$}paku	ra=ara	namana	be=ara	hae	bata=ra	te=o	sinivi\\
     \textsc{art}3.\textsc{sg}	kind	fish	\textsc{rel}	call	1\textsc{pl}.\textsc{in}.\textsc{ipfv}=1\textsc{pl}.\textsc{in}.\textsc{pron}	\textsc{obj}.\textsc{art}2.\textsc{sg}	\textsc{mult}-\textsc{caus}-swing	\textsc{tm}2	\textsc{red}{$\sim$}do	1\textsc{pl}.\textsc{in}.\textsc{ipfv}=1\textsc{pl}.\textsc{in}.\textsc{pron}	ocean	when=1\textsc{pl}.\textsc{in}.\textsc{pron}	board	\textsc{cont}=1\textsc{pl}.\textsc{in}.\textsc{ipfv}	\textsc{prep}=\textsc{art}3.\textsc{sg}	canoe\\
\glt `The kind of fishing that we call \textit{vavaaiku} `letting it swing', we do (it) in the ocean while we are sitting in the canoe.'
\z

\ea\label{ex:text26-2}
Na pakupaku raara bona kanono ae a iri ae sa taba sa ponisi mene dee gunaha bona kanono.\\
\gll na	paku{$\sim$}paku	ra=ara	bona	kanono	ae	a	iri	ae	sa	taba	sa	ponisi	mene	dee	gunaha	bona	kanono\\
     \textsc{tm}2	\textsc{red}{$\sim$}do	1\textsc{pl}.\textsc{in}.\textsc{ipfv}=1\textsc{pl}.\textsc{in}.\textsc{pron}	\textsc{obj}.\textsc{art}2.\textsc{sg}	string	and1	\textsc{art}2.\textsc{sg}	hook	and1	\textsc{part}.\textsc{art}	thing	\textsc{part}.\textsc{art}	heavy	for	carry	go.down	\textsc{obj}.\textsc{art}2.\textsc{sg}	string\\
\glt `We do it with a string and a hook and something heavy to get the string down.'
\z

\ea\label{ex:text26-3}
Beara mei rara tea tahoa, amaa taba vaarii toro vaamatatopo.\\
\gll be=ara	mei	rara	tea	ta-hoa	a=maa	taba	vaarii	toro	vaa-matatopo\\
     when=1\textsc{pl}.\textsc{in}.\textsc{pron}	do.not.yet	1\textsc{pl}.\textsc{in}.\textsc{ipfv}	\textsc{compl}1	\textsc{acaus}-push	\textsc{art}2.\textsc{sg}=\textsc{plm}	thing	\textsc{dem}6.\textsc{pl}	\textsc{tm}7	\textsc{caus}-be.prepared\\
\glt `Before we set off, these things must be prepared.'\footnote{lit. `When we have not set off yet, ...'}
\z

\ea\label{ex:text26-4}
A kanono to vavaaiku ni raara na paku vonaen:\\
\gll a	kanono	to	va-vaa-iku	ni	ra=ara	na	paku	vo=na=en\\
     \textsc{art}2.\textsc{sg}	string	\textsc{rel}	\textsc{mult}-\textsc{caus}-swing	\textsc{appl}	1\textsc{pl}.\textsc{in}.\textsc{ipfv}=1\textsc{pl}.\textsc{in}.\textsc{pron}	\textsc{tm}2	make	like=3\textsc{sg}.\textsc{ipfv}=\textsc{dem}3\\
\glt `The string that we let swing to and fro is made as follows:'
\z

\ea\label{ex:text26-5}
Eara na kisi ni roho rara a ta booru tea kanono, eara repaa kisi mana nia iri putaa.\\
\gll e-ara	na	kisi	ni	roho	rara	a	ta	booru	te=a	kanono	e-ara	re=paa	kisi	mana	ni=a	iri	putaa\\
     \textsc{h}-1\textsc{pl}.\textsc{in}.\textsc{pron}	\textsc{tm}2	tie	\textsc{appl}	before	1\textsc{pl}.\textsc{in}.\textsc{ipfv}	\textsc{art}2.\textsc{sg}	piece	lead	\textsc{prep}=\textsc{art}2.\textsc{sg}	string	\textsc{h}-1\textsc{pl}.\textsc{in}.\textsc{pron}	\textsc{consec}=\textsc{tm}3	tie	eventually	\textsc{appl}=\textsc{art}2.\textsc{sg}	hook	below\\
\glt `We first tie a piece of lead to the string and then we tie the hook below (it).'
\z

\ea\label{ex:text26-6}
Ge beara rake, eara repaa kisi nia iri gina bea ta booru en putaa.\\
\gll ge	be=ara	rake	e-ara	re=paa	kisi	ni=a	iri	gina	be=a	ta	booru	en	putaa\\
     or	when=1\textsc{pl}.\textsc{in}.\textsc{pron}	want	\textsc{h}-1\textsc{pl}.\textsc{in}.\textsc{pron}	\textsc{consec}=\textsc{tm}3	tie	\textsc{appl}=\textsc{art}2.\textsc{sg}	hook	on.top	when=\textsc{art}2.\textsc{sg}	piece	lead	\textsc{dem}3	below\\
\glt `Or if we want, we tie the hook on top while the piece of lead is below (the hook)'
\z

\ea\label{ex:text26-7}
Teo kaku vagana vai, eara na beana rarai.\\
\gll te=o	kaku	vagana	vai	e-ara	na	beana	rara=i\\
     \textsc{prep}=\textsc{art}3.\textsc{sg}	kind	fish	\textsc{dem}	\textsc{h}-1\textsc{pl}.\textsc{in}.\textsc{pron}	\textsc{tm}2	put.bait.on	1\textsc{pl}.\textsc{in}.\textsc{ipfv}=\textsc{dem}1\\
\glt `For this kind of fishing, we bait it (the hook).'
\z

\ea\label{ex:text26-8}
Eara he toro beana nio evo, urita, iobo ge a iana komana vahaa, beara peepee varutarutaa e.\\
\gll e-ara	he	toro	beana	ni=o	urita	iobo	ge	a	iana	komana	vahaa	be=ara	pee{$\sim$}pee	va-ruta{$\sim$}rutaa	e\\
     \textsc{h}-1\textsc{pl}.\textsc{in}.\textsc{pron}	\textsc{advers}	\textsc{tm}7	bait	\textsc{appl}=\textsc{art}3.\textsc{sg}	octopus	sea.sausage	or	\textsc{art}2.\textsc{sg}	fish	even	again	when=1\textsc{pl}.\textsc{in}.\textsc{pron}	\textsc{red}{$\sim$}cut	\textsc{caus}-\textsc{red}{$\sim$}small	3\textsc{sg}.\textsc{pron}\\
\glt `We must bait it with an octopus, a sea-saussage, or even a fish when we have cut them up into little pieces.'
\z

\newpage
\ea\label{ex:text26-9}
Bea maamihu taba vai matatopo, eara repaa hua hiava namana.\\
\gll be=a	maamihu	taba	vai	matatopo	e-ara	re=paa	hua	hiava	namana\\
     when=\textsc{art}2.\textsc{sg}	all.sorts.of	thing	\textsc{dem}6	be.prepared	\textsc{h}-1\textsc{pl}.\textsc{in}.\textsc{pron}	\textsc{consec}=\textsc{tm}3	paddle	go.up	ocean\\
\glt `When everything is prepared, we paddle up to the deep sea.'
\z

\ea\label{ex:text26-10}
Eara repaa vaagunaha a kanono, vaantee koa tobina, eara repaa vaabunobuno koa ni rae, a beana re potee voen na vavaaiku nana.\\
\gll e-ara	re=paa	vaa-gunaha	a	kanono	va-antee	koa	tobina	e-ara	re=paa	vaa-buno{$\sim$}buno	koa	ni	ra=e	a	beana	re	potee	vo=en	na	va{$\sim$}vaa-iku	nana\\
     \textsc{h}-1\textsc{pl}.\textsc{in}.\textsc{pron}	\textsc{consec}=\textsc{tm}3	\textsc{caus}-go.down	\textsc{art}2.\textsc{sg}	string	\textsc{caus}-stop	only	middle	\textsc{h}-1\textsc{pl}.\textsc{in}.\textsc{pron}	\textsc{consec}=\textsc{tm}3	\textsc{caus}-\textsc{red}{$\sim$}move	just	\textsc{appl}	1\textsc{pl}.\textsc{in}.\textsc{ipfv}=3\textsc{sg}.\textsc{pron}	\textsc{art}2.\textsc{sg}	bait	\textsc{consec}	be.like	like=\textsc{dem}3	\textsc{tm}2	\textsc{red}{$\sim$}\textsc{caus}-swing	3\textsc{sg}.\textsc{ipfv}\\
\glt `Then we let the string down, but only half way, and then move it to and fro so that the bait looks like swinging.'
\z

\ea\label{ex:text26-11}
Teebona to kahi avuhu ae tara mau o iana bona, eori repaa maa agomo bona beana.\\
\gll tee-bona	to	kahi	avuhu	ae	tara	ma=u	o	iana	bona	e-ori	re=paa	ma	agomo	bona	beana\\
     \textsc{prep}-\textsc{ana}.\textsc{sg}	\textsc{rel}	\textsc{tm}1	smell	and1	see	hither=\textsc{pfv}	\textsc{art}2.\textsc{pl}	fish	4\textsc{sg}.\textsc{pron}	\textsc{h}-3\textsc{pl}.\textsc{pron}	\textsc{consec}=\textsc{tm}3	come	swallow	\textsc{obj}.\textsc{art}2.\textsc{sg}	bait\\
\glt `So the fish will smell and see it, and then they come and swallow the bait.'
\z

\ea\label{ex:text26-12}
Beara haraa voen na rahi batana bona kanono, eara repaa vaatasu hiava nia kanono,\\
\gll be=ara	haraa	vo=en	na	rahi	bata=na	bona	kanono	e-ara	re=paa	vaatasu	hiava	ni=a	kanono\\
     when=1\textsc{pl}.\textsc{in}.\textsc{pron}	feel	like=\textsc{dem}3	\textsc{tm}2	pull	\textsc{cont}=3\textsc{sg}.\textsc{ipfv}	\textsc{obj}.\textsc{art}2.\textsc{sg}	string	\textsc{h}-1\textsc{pl}.\textsc{in}.\textsc{pron}	\textsc{consec}=\textsc{tm}3	fling	go.up	\textsc{appl}=\textsc{art}2.\textsc{sg}	string\\
\glt `When we feel that (a fish) is pulling the string, we fling the string upwards,'
\z

\newpage
\ea\label{ex:text26-13}
a iri repaa kao hopo rivona iana, eara repaa rahi hiava e kou.\\
\gll a	iri	re=paa	kao	hopo	rivo-n=a	iana	e-ara	re=paa	rahi	hiava	e	kou\\
     \textsc{art}2.\textsc{sg}	hook	\textsc{consec}=\textsc{tm}3	move	enter	mouth-3\textsc{sg}.\textsc{poss}=\textsc{art}2.\textsc{sg}	fish	\textsc{h}-1\textsc{pl}.\textsc{in}.\textsc{pron}	\textsc{consec}=\textsc{tm}3	pull	go.up	3\textsc{sg}.\textsc{pron}	\textsc{emph}2\\
\glt `so that the hook moves into the mouth of the fish, and then we pull it up.'
\z

\ea\label{ex:text26-14}
Eara na antee rara tea paku o kaku vagana vai tea bon, tea roava, teo tarahi ge teo puhana.\\
\gll e-ara	na	antee	rara	tea	paku	o	kaku	vagana	vai	te=a	bon	te=a	roava	te=a	tarahi	ge	te=a	puhana\\
     \textsc{h}-1\textsc{pl}.\textsc{in}.\textsc{pron}	\textsc{tm}2	can	1\textsc{pl}.\textsc{in}.\textsc{ipfv}	\textsc{compl}1	do	\textsc{art}3.\textsc{sg}	kind	fish	\textsc{dem}6	\textsc{prep}=\textsc{art}2.\textsc{sg}	night	\textsc{prep}=\textsc{art}2.\textsc{sg}	daytime	\textsc{prep}=\textsc{art}2.\textsc{sg}	ebbing	or	\textsc{prep}=\textsc{art}2.\textsc{sg}	high.tide\\
\glt `We can do this kind of fishing at night, during the day, when the tide is ebbing or at high tide.'
\z
